\documentclass[11pt,reqno]{amsart}
\usepackage[margin=1in]{geometry}
\usepackage{amsmath,amssymb,amsthm}
\usepackage{graphicx}
\usepackage[colorlinks=true,linkcolor=blue,citecolor=blue,urlcolor=blue]{hyperref}

\theoremstyle{plain}
\newtheorem{theorem}{Theorem}
\newtheorem{proposition}[theorem]{Proposition}
\newtheorem{lemma}[theorem]{Lemma}
\newtheorem{corollary}[theorem]{Corollary}
\newtheorem{conjecture}[theorem]{Conjecture}
\theoremstyle{definition}
\newtheorem{definition}[theorem]{Definition}
\theoremstyle{remark}
\newtheorem{remark}[theorem]{Remark}

\newcommand{\dd}{\,\mathrm{d}}
\newcommand{\R}{\mathbb{R}}
\newcommand{\C}{\mathbb{C}}
\newcommand{\E}{\mathbb{E}}
\newcommand{\Prob}{\mathbb{P}}
\newcommand{\Var}{\operatorname{Var}}
\newcommand{\Ai}{\mathrm{Ai}}
\newcommand{\Bi}{\mathrm{Bi}}
\newcommand{\TW}{\mathrm{TW}}
\newcommand{\W}{\mathcal{W}}
\newcommand{\eps}{\varepsilon}
\newcommand{\Ynode}{Y^\star}

% ---- provenance tags (inherited from the TWbeta paper) ----
\newcommand{\tagPr}{{\footnotesize\textsf{[established here]}}}
\newcommand{\tagP}{{\footnotesize\textsf{[proved elsewhere, cited]}}}
\newcommand{\tagN}{{\footnotesize\textsf{[numerical]}}}
\newcommand{\tagX}{{\footnotesize\textsf{[conjectural / frontier]}}}

\begin{document}

\title[The catastrophe ladder of noise-induced escape]{The catastrophe ladder of
noise-induced escape:\\ from Tracy--Widom to resurgence}
\author{Solomon Williams}
\address{School of Mathematics, University of Edinburgh, Edinburgh EH9 3FD, United Kingdom}
\email{solomonjwilliams635@outlook.com}
\date{Draft, \today}

\begin{abstract}
Noise-induced escape at a slow--fast catastrophe singularity is governed by one reusable
pipeline: geometric blow-up, a Cole--Hopf/Riccati linearization, a stochastic edge operator,
an edge law, and a resurgent trans-series. Instantiated at the fold ($A_2$) the edge law is
exactly Tracy--Widom; at the cusp ($A_3$) it is a new, provably non-integrable law $\W$ whose
only closed form is a resurgent trans-series; and the whole family is organized by the turning
order $q$ into a ladder of \emph{swept multicritical} edge laws with Bessel-$1/(q{+}2)$
skeletons and exact Freidlin--Wentzell tails. We collect the proved spine, the resurgent
structure and its Stokes data, a leading-order edge-variance formula uniform across the ladder,
and a sharp separation between this swept ladder and the genuine (Berry--Upstill)
diffraction-catastrophe ladder. Throughout we distinguish what is proved, what is numerical,
and where the program crosses into a neighbouring field; provenance is tagged on every claim.
\end{abstract}

\maketitle

% =====================================================================
\section{Introduction: one family, one pipeline}\label{sec:intro}

This program occupies one specific intersection: \emph{noise-induced escape at slow--fast
catastrophe singularities}, studied through a single characterization pipeline,
\[
\text{geometric blow-up} \;\to\; \text{Cole--Hopf/Riccati} \;\to\;
\text{stochastic edge operator} \;\to\; \text{edge law} \;\to\; \text{resurgence}.
\]
That pipeline is the program's identity, and this paper is its synthesis. Its purpose is to
state, in one place and with honest provenance, what has been proved, what is established only
numerically, and where the natural questions cross out of the subject into random matrix theory,
resurgence proper, integrable probability, or computational neuroscience.

The organizing parameter is the turning order $q$ of the swept inner equation, indexing the
$A_{q+1}$ singularity: $q{=}1$ fold ($A_2$), $q{=}2$ cusp ($A_3$), $q{=}3$ swallowtail ($A_4$),
and so on. The fold rung is rigorous---the escape law is exactly Tracy--Widom
(\S\ref{sec:fold})---and carries a full resurgent structure with an exact Stokes constant
(\S\ref{sec:twbeta}). The cusp rung is a genuinely new, \emph{provably non-integrable} law $\W$
(\S\ref{sec:cusp}), whose only viable closed form is a resurgent trans-series. Above the cusp the
family is a ladder of swept multicritical edge laws (\S\ref{sec:ladder}); we give a leading-order
edge-variance formula valid on every rung, verify the swallowtail rung, and separate this swept
ladder sharply from the genuine higher diffraction catastrophes (Pearcey, swallowtail function),
which form a distinct family. Section~\ref{sec:physical} anchors the theory in the coupled
neuron model---honestly, as a consolidation within the Berglund--Gentz--Kuehn / Freidlin--Wentzell
framework, not as new phenomena---and \S\ref{sec:open} demarcates the open program.

A word on the two standards of rigour we invoke. ``Proved'' means a theorem, sometimes modulo a
cited black box (Ram\'irez--Rider--Vir\'ag for the stochastic Airy operator; Costin for the
summability of the Painlev\'e~II trans-series). ``Numerical'' means a high-precision computation,
often exact in the sense of Monte-Carlo-free Wiener-chaos evaluation, but not a proof. We never
declare a theorem where a gap remains.

% =====================================================================
\section{The universal pipeline}\label{sec:pipeline}

We state the machinery once, abstractly; each later section is an instantiation.

\smallskip
\noindent\textbf{1. Normal form.} Near a slow--fast fold the reduced dynamics is
$\dd r=(b(\theta)r^2-a(\theta)y)\dd t+\sigma\,\dd W$, $\dd y=-\eps_2 c(\theta)\dd t$, with
$a,b,c>0$ smooth on the slow circle. The deterministic part is the standard folded normal form;
the additive noise is the added ingredient.

\smallskip
\noindent\textbf{2. Geometric blow-up.} The Krupa--Szmolyan quasi-homogeneous blow-up~\cite{KS,JKK}
with weights $(1,2,3)$ on $(r,y,\eps_2)$ resolves the fold; the central rescaling chart $K_2$
($\rho=\eps_2^{1/3}$) carries the inner dynamics.

\smallskip
\noindent\textbf{3. Canonical inner Riccati.} In $K_2$ the SDE becomes
$\dd R=(R^2-Y)\dd T+\eta\,\dd B$, $Y=Y_0-T$, with effective noise $\eta=\sigma/\sqrt{\eps_2}$ (an
exact It\^o change of variables). The deterministic inner solution is the canard
$R=\Ai'(Y)/\Ai(Y)$, escaping at the first Airy zero $a_1\approx-2.3381$.

\smallskip
\noindent\textbf{4. Cole--Hopf linearization.} The substitution $R=-u'/u$ turns the Riccati into a
\emph{linear} stochastic Schr\"odinger equation $u''=(Y-\eta\xi)u$, $\xi=\dd B/\dd T$, with
\emph{no} It\^o--Stratonovich correction; a finite-time blow-up $R\to+\infty$ is exactly a simple
zero (node) of $u$. Escape becomes the first node $\Ynode$ swept downward.

\smallskip
\noindent\textbf{5. Stochastic edge operator.} Writing $x=T$, this is the eigenvalue equation for
the stochastic Airy operator $H_\beta=-\partial_x^2+x+\tfrac{2}{\sqrt\beta}b'(x)$ on $[0,\infty)$
with Dirichlet boundary condition~\cite{RRV}, under the dictionary $\eta=2/\sqrt\beta$.

\smallskip
\noindent\textbf{6. Edge law.} A shooting characterization identifies the \emph{swept} first-node
law with the \emph{fixed} half-line operator's ground state, via the monotone Dirichlet
ground-state function; because the sweep level is deterministic, white-noise stationarity gives
$\Ynode=_d-\Lambda(\beta)=_d\TW_\beta$.

\smallskip
The claim of the pipeline is a universality statement: any slow passage whose inner equation is
the Airy/Riccati canard, perturbed by white noise, Cole--Hopf-linearizes to a Schr\"odinger
operator with white-noise potential, hence to the stochastic Airy operator and the same $\TW_\beta$
exit law, with $\beta$ fixed by the noise-to-curvature ratio. Replacing the fold's linear turning
potential $Y$ by the multicritical $\mathrm{sign}(Y)|Y|^q$ generates the whole ladder of
\S\ref{sec:ladder}.

% =====================================================================
\section{The proved rung: fold $\to$ Tracy--Widom}\label{sec:fold}

\begin{theorem}[Fold escape is Tracy--Widom; RRV cited]\label{thm:fold}
For the inner Riccati with $\eta=2/\sqrt\beta$, the peel-off level satisfies
$\Ynode=_d-\Lambda(\beta)=_d\TW_\beta$, where $\Lambda(\beta)$ is the ground-state eigenvalue of
$H_\beta$. Consequently the early-escape probability is $1-F_\beta(0)$ and, as $\eta\to0$, the
escape rate is $\exp(-c/\eta^2)$ with $c$ a Painlev\'e~II action. \tagP
\end{theorem}

Here $\Lambda(\beta)$ is characterized by Ram\'irez--Rider--Vir\'ag~\cite{RRV}, the leading right-tail
rate agrees with Dumaz--Vir\'ag~\cite{DV}, and at $\beta{=}2$ the law is Tracy--Widom~\cite{TW1994,Bornemann}.

The rate constant $c$ deserves a sharper statement than ``a Painlev\'e~II action, numerically
$5.44$.'' The optimal-fluctuation minimizer solves $-g''+xg=g^3$ on $[0,\infty)$, $g(0)=0$,
$g\to0$, and $c=\int_0^\infty g'^2$.

\begin{proposition}[The rate-constant identity chain]\label{prop:c}
With $g$ the minimizer above,
\[
c=\int_0^\infty g'^2=\int_0^\infty x\,g^2=\tfrac12\int_0^\infty g^4 ,
\]
i.e.\ $A=B=\tfrac12 D$ in the notation $A=\int g'^2$, $B=\int xg^2$, $D=\int g^4$. Numerically
$c=5.4438973582$; $c$ has no elementary closed form (inverse-symbolic search against
$\pi$-powers, $\Gamma(1/3)^3$, $\zeta(3)$, $\log2$ to $10$ digits is null). \tagPr
\end{proposition}

\begin{proof}
Multiplying the Euler--Lagrange equation by $g$ and integrating by parts gives $A+B=D$ (identity I).
Multiplying by $xg'$ and integrating by parts gives $\tfrac12A+\tfrac14D=B$ (identity II, the
Pohozaev/dilation identity). Eliminating $D=A+B$ from (II) yields $\tfrac34A=\tfrac34B$, so $A=B$
and $D=2A$. \qed
\end{proof}

At the process level, the leading escape of a spatially extended folded medium realizes the full
Airy$_2$ process---not merely its $\TW_2$ marginal---\emph{if and only if} the transverse coupling
is disordered, long-range, and stochastically drifting; the increment-variance exponent (Brownian
$1$ vs.\ ballistic $2$), not the shared marginal, is the discriminant. These process-level
statements are numerical; rigorous convergence to the Airy$_2$ process (Brownian Gibbs / line
ensemble) is open (\S\ref{sec:open}).

% =====================================================================
\section{Resurgent depth at the fold}\label{sec:twbeta}

The $\TW_\beta$ right tail carries resurgent structure in \emph{two} independent directions: a
\emph{level} direction ($a\to\infty$ at fixed $\beta$), governed by Painlev\'e~II and closed by
Costin's theorem~\cite{Costin,CCK}, and a \emph{noise} direction ($\beta\to\infty$ at fixed $a$),
whose $1/\beta$ loop expansion is separately median Borel summable by elementary
means~\cite{Sokal}. The all-orders soft-tail form is that of Borot--Nadal~\cite{BorotNadal}, with the
large-$\beta$ reduction to PII from Comtet--Le~Doussal--Smith~\cite{CLDS}. The level direction rests
on an exact linearization: the Jacobi operator $\mathcal J=-\partial_x^2+6p_\star^2-2(x+a)$
\emph{is} the Painlev\'e~II linearization $\mu^2(-\partial_s^2+6q^2+s)$, the instanton being the
$\alpha=\tfrac12$ Airy special solution of PII~\cite{FN1980,JMU1981,FIKN,Clarkson}.

\begin{theorem}[Stokes constant of the backbone; exact]\label{thm:C}
The $s\to-\infty$ Hastings--McLeod coefficients are non-sign-alternating; the Borel singularity
sits on the positive axis at action $A=2\sqrt2/3$, index $\gamma=-\tfrac12$, with Stokes constant
\[
C=\sqrt{\tfrac{2}{3\pi^3}}=\tfrac1\pi\sqrt{\tfrac{2}{3\pi}}=0.146632271193848478\ldots
\]
matching the Hastings--McLeod value~\cite{CleriDunne,Dunne2025,KapaevHM} and refuting the natural
guess $1/(4\sqrt\pi)$. The $a$-dependence is algebraic: $C$ factorizes as a transverse Kramers
determinant $1/\pi$ times a Gelfand--Yaglom zero-mode Jacobian $\sqrt{2/(3\pi)}=\|\psi_0\|/\sqrt{2\pi}$
of the $\tanh$-kink escape instanton~\cite{SGG}. \tagPr\,/\,\tagP
\end{theorem}

The $O(1/\beta)$ noise-averaged shift of the connection root admits a closed \emph{form} (though
not a closed number).

\begin{proposition}[The connection-root shift $\Omega_1$]\label{prop:om1}
Writing $\E[\Lambda_0]=\Lambda_0^{\det}+\beta^{-1}\Omega_1+O(\beta^{-2})$ with
$\Lambda_0=2.3381074105$, second-order Rayleigh--Schr\"odinger perturbation over the Airy
eigenbasis gives $\Omega_1=-4\int_0^\infty\psi_0^2\,G_{\mathrm{red}}$, which reduces to
\[
\Omega_1=\frac{4\pi}{\Ai'(-\Lambda_0)^2}
\left[\frac{\Bi'(-\Lambda_0)}{\Ai'(-\Lambda_0)}\int_{a_1}^\infty\!\Ai^4
-\int_{a_1}^\infty\!\Ai^3\Bi\right],
\]
the third natural term $\int_{a_1}^\infty\Ai^3\Ai'$ vanishing exactly because $a_1$ is a zero of
$\Ai$. Numerically $\Omega_1=-1.124813904375434066$ ($27$ digits). The quartic moment
$\int_{a_1}^\infty\Ai^4$ is provably not boundary-reducible, so $\Omega_1$ is a genuine
Airy-quartic transcendental. \tagPr
\end{proposition}

\begin{proposition}[Which Painlev\'e~II: saddle vs.\ distribution]\label{prop:alpha}
The leading tail \emph{saddle} is the $\alpha=\tfrac12$ Airy special solution; the full $\beta{=}2$
\emph{distribution} is the $\alpha=0$ Hastings--McLeod transcendent. These are consistent: the
$\alpha=\tfrac12$ saddle action gives
$\mathcal T_\beta(a)\sim a^{-3\beta/4}e^{-(2\beta/3)a^{3/2}}$, which at $\beta{=}2$ reproduces the
classical $\TW_2$ right tail $e^{-(4/3)a^{3/2}}/(16\pi a^{3/2})$ in both exponent and prefactor
power (verified: the ratio $\to1$, e.g.\ $0.941,0.957,0.967$ at $a=8,10,12$). Thus $\alpha=\tfrac12$
governs the saddle/rate and $\alpha=0$ the distribution; this is the ordinary saddle-vs-distribution
distinction, not a $\beta$-dependence. \tagPr
\end{proposition}

The two directions share the single action $\Phi(a)=\tfrac23a^{3/2}$, and the tail is a
Borel-summable resurgent trans-series, uniformly in $a$ for $a\ge a_0\approx1$.

% =====================================================================
\section{The cusp: a provably non-integrable law $\W$}\label{sec:cusp}

At the cusp the inner equation is the stochastic Weber operator
$u''=(\mathrm{sign}(Y)Y^2-\eta\xi)u$, $\beta=4/\eta^2$, and the escape law $\W_\beta$ is a genuinely
new object. Two obstructions make this a pair of theorems, not a gap.

\begin{proposition}[No integrable closed form]\label{prop:noint}
\emph{(i)} The deterministic Weber skeleton is an isomonodromy fixed point~\cite{JMU1981} (the only
relevant deformation flows cusp$\to$fold), so $\W$ is not the $\sigma$-function of a Painlev\'e
transcendent.
\emph{(ii)} $\W$ is a first-passage law of an operator unbounded below (escape occurs down the
$-Y^2$ side), hence not a soft-edge Fredholm determinant. \tagP
\end{proposition}

In particular the ladder analogy ``$\mathrm{Airy}{:}\mathrm{PII}::\mathrm{Weber}{:}\mathrm{PIV}$''
does \emph{not} extend to the law: $\W$ is not a Painlev\'e~IV $\sigma$-function. What survives is a
closed-form \emph{deterministic connection}. The recessive log-derivative at $Y=0$ on the confluent
($Y>0$) side is $R(\lambda)=-2\,\Gamma(\tfrac34-\tfrac\lambda4)/\Gamma(\tfrac14-\tfrac\lambda4)$; the
oscillatory ($Y<0$) side is the confluent equation rotated by $e^{-i\pi/4}$ with $\lambda\to i\lambda$,
giving $L(\lambda)=2e^{-i\pi/4}\,\Gamma(\tfrac34-\tfrac{i\lambda}4)/\Gamma(\tfrac14-\tfrac{i\lambda}4)$.

\begin{proposition}[The cusp resonance root is an exact Gamma-equation root]\label{prop:lz}
The resonance condition $L(\lambda)=R(\lambda)$ is the explicit equation
\[
e^{-i\pi/4}\,\frac{\Gamma(\tfrac34-\tfrac{i\lambda}4)}{\Gamma(\tfrac14-\tfrac{i\lambda}4)}
+\frac{\Gamma(\tfrac34-\tfrac\lambda4)}{\Gamma(\tfrac14-\tfrac\lambda4)}=0,
\]
with root $\lambda_0=0.8896052163757951\,(1-i)$. Its phase is \emph{exactly} $-\pi/4$: on the ray
$\lambda=\mu(1-i)$ the two $\Gamma$-arguments are complex conjugates, so the equation reduces to the
real condition $\arg\bigl[\Gamma(\tfrac34-\tfrac{\mu(1-i)}4)/\Gamma(\tfrac14-\tfrac{\mu(1-i)}4)\bigr]
=3\pi/8$. Moreover $\arg D'(\lambda_0)=-7\pi/8$ exactly. The modulus $|\lambda_0|=1.2580918$ has no
elementary closed form. \tagPr
\end{proposition}

The stochastic data are then a resurgent trans-series in $\eta^2$. The leading weak-noise variance,
$\Var(\Ynode)=\eta^2(v_0+v_1\eta^2+\cdots)$, has an exact one-loop form (see \S\ref{sec:ladder});
$v_0=0.13429234295678519$, transcendental. The first noise correction to the connection root,
$\E[\lambda]=\lambda_0+\eta^2\Omega_2+O(\eta^4)$, is $\Omega_2=-0.451+0.352i$ (also transcendental);
Proposition~\ref{prop:lz} makes its prefactor $-1/D'(\lambda_0)$ exact (the deterministic Weber
Stokes normalization is that of Hao~\cite{Hao}). A median Borel resummation of seven exact weak-noise
coefficients reconstructs $\Var$ at $\beta{=}2$---outside the perturbative radius---to about $1\%$,
which is the sense in which the trans-series \emph{is} $\W$. \tagN

% =====================================================================
\section{The ladder and its two directions}\label{sec:ladder}

Replacing the turning potential by $V_q(Y)=\mathrm{sign}(Y)|Y|^q$ gives the swept multicritical
family $u''=(V_q-\eta\xi)u$. Three structural layers are uniform across it:
\emph{(1)} the deterministic skeleton is a Bessel function of order $1/(q{+}2)$, and the peel-off
levels are its zeros (at $q{=}1$, the exact Airy zeros); \emph{(2)} the exact Freidlin--Wentzell
tails are
\[
\Prob(Y^\star<-s)\approx\exp\Bigl(-\frac{|s|^{2q+1}}{2(2q{+}1)\,\eta^2}\Bigr)
=\exp\Bigl(-\frac{\beta\,|s|^{2q+1}}{8(2q{+}1)}\Bigr)\quad\text{(left)},
\qquad \exp(-\tfrac43|s|^{3q/2})\quad\text{(right)};
\]
\emph{(3)} the edge law is a resurgent trans-series. To these we add a uniform leading-order
statement.

\begin{remark}[The left-tail constant, made $\eta$-explicit]\label{rem:fwconstant}
Earlier drafts of this section quoted the left tail as $\exp(-|s|^{2q+1}/[4(2q{+}1)])$, with no
$\eta$. That expression is the $\beta{=}2$ specialisation of the law above, not a general-$\eta$
statement, and taken literally it disagrees with the cusp paper's
$-\log\Prob\to s^5/(10\eta^2)=\beta s^5/40$. The $\eta$-explicit form reconciles the two: at $q=2$
it gives exactly $\beta s^5/40$, and at $\beta=2$ it reproduces the earlier numbers. The constant
$1/[2(2q{+}1)]$ is verified numerically for $q=1,\dots,5$ by boundary-value solution of the
Euler--Lagrange system (measured/predicted ratios $0.839$, $0.930$, $0.981$, $0.997$, $0.9995$,
approaching $1$ from below since at fixed $s$ the steeper $q$ reaches its asymptote sooner).
\tagN
\end{remark}

\begin{proposition}[Uniform leading edge-variance]\label{prop:v0q}
Let $\varphi_q$ be the recessive solution of $\varphi''=\mathrm{sign}(Y)|Y|^q\varphi$ and $Y_\star^0$
its first node swept from $+\infty$. Then the leading weak-noise variance coefficient is
\[
v_0^{(q)}=\frac{\int_{Y_\star^0}^\infty \varphi_q(Y)^4\,\dd Y}{\varphi_q'(Y_\star^0)^4}.
\]
The derivation is a first-order node-shift computation in which the Wronskian identity at the node,
$W=-\varphi'(Y_\star^0)\psi(Y_\star^0)$, cancels the dominant-solution factor. Numerically
$v_0^{(1)}=0.41743$, $v_0^{(2)}=0.13429$, $v_0^{(3)}=0.04953$, $v_0^{(4)}=0.02097$; the $q{=}1$
node is the exact first Airy zero. \tagPr
\end{proposition}

The swallowtail rung ($q{=}3$, $A_4$) is verified in full: the skeleton is Bessel-$1/5$ (its first
five peel-off levels match the matched Bessel combination---the symmetric $J_{1/5}+J_{-1/5}$---to
$10^{-14}$), and the standardized escape moments at $\beta{=}2$ are skewness $+0.951$ (predicted
$+0.95$) and excess kurtosis $+0.156$ (predicted $+0.16$), from $2.8\times10^6$ samples. \tagN

\begin{figure}[t]
\centering
\includegraphics[width=\textwidth]{cap_fig_skeleton.pdf}
\caption{The swept ladder. (A) The swallowtail ($q{=}3$) deterministic skeleton
$\varphi''=Y^3\varphi$: recessive (decaying) for $Y>0$, oscillatory for $Y<0$; its peel-off levels
are the zeros of the Bessel-$1/5$ combination (first five agree to $10^{-14}$). (B) The uniform
leading edge-variance $v_0^{(q)}$ of Proposition~\ref{prop:v0q} across the ladder.}
\label{fig:skeleton}
\end{figure}

\subsection*{Swept versus genuine.} This swept ladder must not be confused with the genuine
Berry--Upstill diffraction-catastrophe ladder~\cite{BerryUpstill} $u^{(n)}=(Y-\eta\xi)u$, an $n$-th
order equation:
fold ($A_2$) is Airy, cusp ($A_3$) is \emph{Pearcey} (third order), swallowtail ($A_4$) is the
swallowtail function (fourth order). The two families \emph{coincide only at the fold and diverge in
opposite directions}. As $q$ rises the swept laws grow more skewed and sub-Gaussian (skewness
$+0.20,+0.61,+0.95$); as the order $n$ rises the genuine laws trend toward Gaussian (skewness
$+0.225,+0.120,+0.058$, the first matching $\TW_2$ and validating the integrator). \tagN\
The interpretation is structural: the scalar one-fast-variable pipeline of \S\ref{sec:pipeline}
(Cole--Hopf of a scalar Riccati) produces a \emph{second-order} Schr\"odinger operator, hence the
swept ladder; the genuine swallowtail requires a fourth-order linear inner equation, which does not
arise from a scalar canard and needs higher-order fast structure. The family therefore does not
terminate---it \emph{splits}, and genuine $A_4$ is the point at which the scalar theory bridges out
(Figure~\ref{fig:ladders}).

\begin{figure}[t]
\centering
\includegraphics[width=\textwidth]{cap_fig_ladders.pdf}
\caption{The two ladders. (A) Peel-off skewness: the swept turning-point ladder (2nd order) grows
more sub-Gaussian with $q$, while the genuine Berry--Upstill ladder ($n$-th order) trends toward
Gaussian; they coincide only at the fold ($=\TW_2$). (B) The swept standardized edge laws
fold$\to$cusp$\to$swallowtail, from $2.8\times10^6$ Monte-Carlo escapes each at $\beta{=}2$.}
\label{fig:ladders}
\end{figure}

% =====================================================================
\section{Physical realization, honestly positioned}\label{sec:physical}

The turning order is not abstract: in the coupled fold--Hopf--relaxation neuron model the electrical
coupling $g$ sets the turning structure through $\Delta(g)=2\sqrt{-2g/3}$, so tuning $g$ moves the
system along the swept ladder---fold ($\TW$) escape statistics at weak coupling, through the cusp
(sub-Gaussian, negative excess kurtosis) at $g_{\mathrm{crit}}$. The sign of the escape-timing
excess kurtosis is the observable rung signature.

We position this carefully. The single most novel content of the program is the
\emph{theory spine}: the fold$=\TW$ theorem, the resurgent structure with its exact Stokes constant,
and the non-integrable cusp law $\W$. The \emph{neuron realization}, by contrast, sits squarely
within the Berglund--Gentz~\cite{BG} and Berglund--Gentz--Kuehn~\cite{BGK} slow--fast noise
framework; the folded-node / mixed-mode noise phenomenology it touches is their territory, and the
clean closed forms found there (e.g.\ the universal fold-escape prefactor $\sqrt{4\pi\ln2}$) are
consolidations within a known framework, not new phenomena. We claim the demonstration---that the
ladder \emph{lives} in the neurons---not novelty in the neuronal modelling.

% =====================================================================
\section{The demarcated open program}\label{sec:open}

What closes the book is the account above: the proved spine, the resurgent structure to the depth the
present methods reach, the ladder and its separation from the genuine catastrophes, and the physical
anchor. What remains is an \emph{explicitly bounded} research program, and every substantial direction
in it crosses into a different field.

\begin{itemize}
\item \emph{Random matrices / topological recursion.} The all-genus coefficients $R_m$ of the
Borot--Nadal soft-tail formula~\cite{BorotNadal} require the full topological-recursion machinery;
this is unreachable from within the present methods.
\item \emph{Resurgence proper.} A direct exact-WKB / Voros-symbol resurgence for the random operator
itself would treat the white-noise potential without averaging; no exact-WKB framework exists for a
nowhere-differentiable potential (the Voros symbols $\int\sqrt Q$ are ill-defined pathwise). The
scoped attack is parametric-Stokes theory with the noise coupling as deformation parameter.
\item \emph{Integrable probability / KPZ.} Process-level convergence to the Airy$_2$ (and,
conjecturally, Pearcey) process needs a Brownian-Gibbs / line-ensemble property for the
folded-medium edge; only the marginals and the increment-variance exponent are in hand.
\item \emph{Computational neuroscience.} The realization in actual networks and data.
\end{itemize}

Two residuals are softer and internal: the rough-path rigidity of the rough-potential isomonodromy
(the $O(\eps^2)$ location shift is $\delta(0)$-divergent but has a finite Wick-ordered form), and the
nonlinear first-passage sensitivity bound $R_{\mathrm{nl}}$ that stands between the proved
sub-Gaussian inner moments and a fully closed connection-phase cumulant bound. Neither is a new
framework; both are writeable with effort. Everything past them is a deliberate crossing, chosen on
its own merits, not a deficiency of the account here.

\begin{thebibliography}{99}

\bibitem{KS} M.~Krupa, P.~Szmolyan, \emph{Extending geometric singular perturbation theory to
nonhyperbolic points---fold and canard points in two dimensions}, SIAM J. Math. Anal. \textbf{33}
(2001) 286--314.

\bibitem{JKK} S.~Jelbart, C.~Kuehn, N.~Kuntz, \emph{Geometric blow-up of a dynamic
Hopf/fold-of-cycles}, arXiv:2208.01361 (2024).

\bibitem{RRV} J.~Ram\'irez, B.~Rider, B.~Vir\'ag, \emph{Beta ensembles, stochastic Airy spectrum,
and a diffusion}, J. Amer. Math. Soc. \textbf{24} (2011) 919--944; arXiv:math/0607331.

\bibitem{DV} L.~Dumaz, B.~Vir\'ag, \emph{The right tail exponent of the Tracy--Widom-$\beta$
distribution}, Ann. Inst. H. Poincar\'e Probab. Statist. \textbf{49} (2013) 915--933; arXiv:1102.4818.

\bibitem{TW1994} C.~A.~Tracy, H.~Widom, \emph{Level-spacing distributions and the Airy kernel},
Comm. Math. Phys. \textbf{159} (1994) 151--174.

\bibitem{Bornemann} F.~Bornemann, \emph{On the numerical evaluation of distributions in random
matrix theory: a review}, Markov Process. Related Fields \textbf{16} (2010) 803--866.

\bibitem{BorotNadal} G.~Borot, C.~Nadal, \emph{Right tail asymptotic expansion of Tracy--Widom beta
laws}, Random Matrices Theory Appl. \textbf{1} (2012) 1250006; arXiv:1111.2761.

\bibitem{CLDS} A.~Comtet, P.~Le Doussal, N.~R.~Smith, \emph{The Tracy--Widom distribution at large
Dyson index}, J. Stat. Phys. \textbf{193} (2026) Art.~46; arXiv:2510.14433.

\bibitem{Costin} O.~Costin, \emph{On Borel summation and Stokes phenomena of nonlinear differential
systems}, Duke Math. J. \textbf{93} (1998) 289--344; arXiv:math/0608408.

\bibitem{CCK} O.~Costin, R.~D.~Costin, M.~Kohut, \emph{Rigorous bounds of Stokes constants for some
nonlinear ODEs at rank-one irregular singularities}, Proc. R. Soc. A \textbf{460} (2004) 3631--3641;
arXiv:math/0608316.

\bibitem{Sokal} A.~D.~Sokal, \emph{An improvement of Watson's theorem on Borel summability},
J. Math. Phys. \textbf{21} (1980) 261--263.

\bibitem{CleriDunne} N.~J.~Cleri, G.~V.~Dunne, \emph{Resurgent trans-series for generalized
Hastings--McLeod solutions}, J. Phys. A: Math. Theor. \textbf{53} (2020) 305201; arXiv:2002.06270.

\bibitem{Dunne2025} G.~V.~Dunne, \emph{Introductory lectures on resurgence}, arXiv:2511.15528 (2025).

\bibitem{KapaevHM} A.~A.~Kapaev, \emph{Quasi-linear Stokes phenomenon for the Hastings--McLeod
solution of the second Painlev\'e equation}, arXiv:nlin/0411009 (2004).

\bibitem{SGG} T.~Schorlepp, T.~Grafke, R.~Grauer, \emph{Symmetries and zero modes in sample path
large deviations}, J. Stat. Phys. \textbf{190} (2023) 50; arXiv:2208.08413.

\bibitem{FN1980} H.~Flaschka, A.~C.~Newell, \emph{Monodromy- and spectrum-preserving deformations. I},
Comm. Math. Phys. \textbf{76} (1980) 65--116.

\bibitem{JMU1981} M.~Jimbo, T.~Miwa, K.~Ueno, \emph{Monodromy preserving deformation of linear
ordinary differential equations with rational coefficients. I}, Physica D \textbf{2} (1981) 306--352.

\bibitem{FIKN} A.~S.~Fokas, A.~R.~Its, A.~A.~Kapaev, V.~Yu.~Novokshenov, \emph{Painlev\'e
Transcendents: The Riemann--Hilbert Approach}, Math. Surveys Monogr. \textbf{128}, AMS, 2006.

\bibitem{Clarkson} P.~A.~Clarkson, \emph{Painlev\'e equations---nonlinear special functions},
J. Comput. Appl. Math. \textbf{153} (2003) 127--140.

\bibitem{Hao} Q.~Hao, \emph{Exact WKB of solutions by Borel summation and open TBA},
arXiv:2507.06922 (2025).

\bibitem{BerryUpstill} M.~V.~Berry, C.~Upstill, \emph{Catastrophe optics: morphologies of caustics
and their diffraction patterns}, Prog. Optics \textbf{18} (1980) 257--346.

\bibitem{BG} N.~Berglund, B.~Gentz, \emph{Noise-Induced Phenomena in Slow--Fast Dynamical Systems:
A Sample-Paths Approach}, Springer, 2006.

\bibitem{BGK} N.~Berglund, B.~Gentz, C.~Kuehn, \emph{Hunting French ducks in a noisy environment},
J. Differential Equations \textbf{252} (2012) 4786--4841; \emph{From random Poincar\'e maps to
stochastic mixed-mode-oscillation patterns}, J. Dynam. Differential Equations \textbf{27} (2015)
83--136.

\end{thebibliography}

\end{document}
